La Ley General de Masas queda formulada como una transformación genealógica y
espectral, no como una colección de números. Su dominio efectivo es la cadena
\[
\text{APP}\longrightarrow\text{TRIT}\longrightarrow\text{TPK}
\longrightarrow\Gamma^{\rm surv}\longrightarrow\Gamma^{\rm obs}
\longrightarrow L_{108}\longrightarrow\sigma\in\mathbb Z^5
\longrightarrow\chi_{\rm full}\longrightarrow\widehat M,
\]
seguida, cuando la representación física lo permite, por el funcional que
extrae una coordenada escalar de la recta de masa. Cada flecha conserva su
dominio, su orientación, su memoria y su estatuto; ninguna cifra externa elige
retroactivamente una ruta, una firma o un coeficiente armónico.

El atlas finito contiene 81 celdas, 56 familias de ruta, 13 multisecciones y
324 rutas orientadas. El punto central electrónico es la única fibra puntual
equivariante. Las demás clases elementales se realizan naturalmente como
fibras, órbitas o multisecciones, de modo que su salida primaria puede ser un
operador o un espectro antes de ser una masa escalar. Esta diferencia no es una
carencia de dominio, sino una propiedad de la geometría interna.

En la fibra central, el selector basal y el lector bidireccional producen
\[
\mathfrak e_0
=\frac{\sqrt3}{4}
 \left[\exp\!\left(\frac{\varphi}{\pi^2}\right)
       -22\alpha_{\rm HMT}^3\right](1+15\Delta_4),
\qquad
\kappa_e=80-54\alpha_{\rm HMT}+6\Delta_4,
\]
\[
\mathfrak e_e^\beta=\mathfrak e_0R_{\rm act}^{\kappa_e},
\qquad
\varsigma_{u_E}(\mathfrak e_e^\beta)
=0.510998950690000808608\ldots\ \mathrm{MeV}.
\]
La escala absoluta de acción permanece en la rama
\(\Sigma_H=\varphi\alpha_{\rm HMT}^{16}\), de orden decimal \(-34\), y se
acopla al reloj mediante \(\hbar_{\rm int}\omega\); al formar cocientes de
masas, toda sección común cancela y sólo sobreviven los caracteres, los
lectores y los refinamientos relativos correctamente tipados.

La jerarquía armónica no consta de ocho correcciones aisladas. Es el semigrupo
global
\[
\langle90,120\rangle
=\{90,120\}\cup\{30r:r\ge6\},
\qquad
c_n=q_-^n+q_+^n,
\quad s_n=q_-^n-q_+^n,
\]
actuando sobre un único operador bicapa. Los niveles visibles son testigos de
esa jerarquía; su completitud representacional no sustituye al mapa físico que
debe transportar ruta, acarreo, memoria y devanado a los coeficientes de cada
estado.

Las realizaciones físicas tipadas, las diecisiete evaluaciones condicionadas, las
dos secciones nulas y los sectores neutrínicos, compuestos, resonantes,
topológicos y virtuales quedan así reunidos sin aplanar sus codominios. El
inventario posterior de 471 masas y 384 anchuras conserva el contraste
estado por estado, pero no se confunde con un censo de predicciones. Una
comparación cuantitativa adquiere fuerza física sólo después de declarar
observable, unidad, esquema, escala, incertidumbre, covarianza y regla de
realización.

La conclusión estructural es precisa: la masa es la medida espectral de una
persistencia bajo acoplamiento. La partícula es la sección persistente; la ruta,
su historia observable; la firma, su lectura entera; el carácter, su acción
adimensional; las torres, su resolución armónica; y la carta dimensional, la
realización de esa genealogía en unidades físicas. El atlas y los inventarios
precedentes permiten reproducir cada afirmación sin reducir la teoría a una
tabla abreviada ni convertir una interfaz abierta en un teorema cerrado.
